La detección temprana de fallas en correas transportadoras es fundamental para evitar paradas no programadas y pérdidas operativas en industrias críticas. Los métodos tradicionales basados en inspección visual presentan limitaciones en velocidad, costo y precisión. En contraste, las técnicas de aprendizaje profundo, especialmente aquellas basadas en redes neuronales convolucionales (CNN) y arquitecturas \textit{Transformer}, ofrecen soluciones automatizadas y escalables para la clasificación de fallas mediante visión computacional, favoreciendo el mantenimiento predictivo y la supervisión continua.\\

Este trabajo de memoria de título se enfoca en implementar modelos de detección de fallas en correas transportadoras mediante técnicas de visión computacional y aprendizaje profundo. El objetivo general es implementar y evaluar algoritmos de detección de objetos basados en inteligencia artificial. Entre los objetivos específicos destacan la implementación y entrenamiento con bases de datos etiquetadas, la programación de modelos basados en CNN (YOLO y SSD) y modelos \textit{Transformer} (DETR y DINO). Por último, se realiza una evaluación comparativa mediante métricas estándar de visión por computadora.\\

La metodología contempla la implementación y el entrenamiento de modelos en entornos de cómputo especializados, utilizando bases de datos segmentadas rigurosamente en entrenamiento, validación y prueba. Los resultados posicionan al modelo YOLOv5 como la arquitectura más eficiente y de mayor viabilidad industrial, superando a SSD512 y a los modelos de atención global. Si bien los detectores basados en \textit{Transformers} demostraron ventajas teóricas al aislar defectos complejos como perforaciones, su desempeño general se vio penalizado por la escala reducida de los datos y las restricciones de capacidad física (\textit{hardware}). A partir de estos hallazgos, se propone como trabajo futuro la exploración de arquitecturas híbridas y una optimización de hiperparámetros. Para garantizar la reproducibilidad, la estructura de trabajo, registros de métricas y detalles se encuentran almacenados en un repositorio de código abierto (\href{https://github.com/FernandoN23/Implementation-of-Object-Detection-for-Conveyor-Belt-Defect-Detection}{\textit{Implementation of Object Detection for Conveyor Belt Defect Detection}}).